\section{The Symbols Nobody Provides}
\label{sec:ble-symbols}

The bindings compile, the blobs are in the link path, and the linker still fails — on functions that
exist nowhere in the SDK's public surface. This section is the complete answer for this board, because
it cost a full session to assemble.

\subsection{The CMSIS Helpers}

\begin{minted}{text}
undefined symbol: CORE_EnterAtomic
  >>> referenced by bgbuf.c.obj:(bgbuf_free)
undefined symbol: CORE_ExitAtomic
  >>> referenced by ll_hci.c.obj:(ll_hciDataFromHost)
\end{minted}

In a vendor project these are \codeinline{static inline} in the CMSIS headers, so the C compiler
inlines them and no symbol is ever needed. The blob was compiled against those headers, so its object
code contains \emph{references}, and the application is expected to provide the definitions.

The definitions are four lines each, and the two that matter differ:

\begin{minted}{c}
uint32_t CORE_EnterCritical(void) { /* masks with BASEPRI */ }
uint32_t CORE_EnterAtomic(void)   { /* masks with PRIMASK */ }
\end{minted}

\textbf{This difference is not cosmetic.} \codeinline{BASEPRI} masks interrupts at or below a chosen
priority and leaves priority~0 running; \codeinline{PRIMASK} masks everything. A stack that enters a
buffer-manipulating critical section with the wrong one still links, still works at idle, and corrupts
its attribute buffer once a connection is active. Implement them from the CMSIS semantics, not from a
guess about what ``critical'' means.

\subsection{Provided Once, for Every Stack}

This board links two blob families — Bluetooth and RAIL — and both want the same helpers. There are
two ways to get that wrong:

\begin{itemize}
    \item \textbf{Compile the C file in each crate.} Two archives with the same symbols: either a
          duplicate-symbol link error, or one silently wins and the other's semantics are lost.
    \item \textbf{Give each its own copy.} Works, until the two copies drift and only one is updated.
\end{itemize}

This project's answer is a build-support crate that compiles the file once, which both stacks depend
on. A crate with no Rust API, whose entire contract is ``linking me gives you these symbols'':

\begin{minted}{toml}
[dependencies]
cmsis-stubs = { path = "../cmsis-stubs" }
\end{minted}

\subsection{The \#[link] Directive That Is Not Obvious}

Having the crate as a dependency is not enough: nothing in Rust \emph{calls} those symbols, so the
linker never pulls the archive in, and every symbol comes back undefined while the archive sits in the
search path. The missing piece is an explicit directive:

\begin{minted}{rust}
#[link(name = "cmsis_stubs", kind = "static")]
extern "C" {}
\end{minted}

Five lines that do nothing at runtime and everything at link time. They belong in the crate whose
objects need the symbols — the bindings crate — not in the helper crate.

\subsection{The Rule That Generalises}

\begin{enumerate}
    \item Read the link errors in full, once, and list every symbol. The linker names the object file
          that wanted each one (\codeinline{bgbuf.c.obj}, \codeinline{ll_hci.c.obj}), which tells you
          which layer is asking.
    \item Decide, per symbol, whether it is \emph{data you must supply} (CMSIS, emlib) or \emph{a
          library you must vendor} (PSA crypto). They have completely different fixes.
    \item Supply data once for the whole binary; vendor libraries with their licence in mind.
    \item Declare the link directive where the reference is, and comment why — because the next reader
          will see an empty \codeinline{extern} block and delete it.
\end{enumerate}
